% Solutions/hints for ch14 exercises (assembled into Appendix C)
\begin{solution}{exr:ch14-braking}
(a) $s_b = v^2/(2a_b) = 4/4 = 1$~m. (b) The command is held for $\dt$ before braking starts: $v\dt + v^2/(2a_b) = 0.5 + 1 = 1.5$~m. (c) Continuous condition: $\sqrt{2 \cdot 0.608 \cdot 2} = 1.56$~m/s; discrete condition: $-0.5 + \sqrt{0.25 + 4 \cdot 0.608} = 1.14$~m/s. (d) The physics is the same, but in a digital loop the drone cannot begin to brake before the current command has run its course, so the reaction delay $\dt$ costs $v\dt$ of the free distance.
\end{solution}
\begin{solution}{exr:ch14-scores}
With $(0.8, 0.1, 0.1)$: $G(1.0, 0) = 0.800 + 0.020 + 0.050 = 0.870$, $G(1.0, 0.5) = 0.675 + 0.100 + 0.056 = 0.831$, $G(0.75, 0.5) = 0.644 + 0.100 + 0.045 = 0.789$. The straight command wins (it is also the winner among all 16 admissible candidates): the heading-dominant drone keeps flying at the obstacle and swerves only when the braking test forces it to. With $(0.2, 0.6, 0.2)$ the scores are $0.422$, $0.881$ and $0.851$: the swerve wins because clearance now counts three times as much as heading. The straight-flying drone does not crash because $(1.0, 0)$ is admissible, its braking distance of $0.5$~m fits into the $0.608$~m of free distance, and by the safety theorem the braking candidate remains admissible in every later step.
\end{solution}
\begin{solution}{exr:ch14-hysteresis}
(a) The scene is mirror-symmetric about the goal line, so the free distance, the heading angle and the speed of $(v_x, v_y)$ and $(v_x, -v_y)$ coincide and the tie is broken by the order of the grid. After one step the drone is off the axis by $v_y\dt$; the side it moved to now has a slightly worse heading and a slightly better clearance, and depending on the weights the mirror candidate can win the next step and move the drone back. (b) If $\abs{G(\vel_1) - G(\vel_2)} < \eta$ and $\vel_a = \vel_1$ is admissible, then $G(\vel_a) \ge G(\vel_2) - \eta$ and the rule keeps $\vel_1$; the alternation can only start once the scores differ by at least $\eta$. (c) Hysteresis only decides between candidates whose scores are within $\eta$ of each other; in the U-trap the commands that would leave the U score far below the commands along the wall for every $\eta$, so the drone keeps shuttling or settles near an arm. With \code{DwaParams(hysteresis=0.05)} the run on \code{utrap\_scene()} still ends inside the U after 300 steps, having flown about 36~m instead of 80~m: the oscillation is damped, the trap remains.
\end{solution}
